% ===================== page 1 =====================
\printedpage{--} % no page number printed (title page)
% head: (none)

\begin{center}
{\Large\textsf{THE DERIVATORS}}

\bigskip

\textbf{ALEXANDRE GROTHENDIECK}

\bigskip\bigskip

\textbf{Chapter I}

\medskip

\textbf{Generalities on (pre)derivators}
\end{center}

\clearpage

% ===================== page 2 =====================
\printedpage{--} % no page number printed (editors' note)
% head: (none)

This text was deciphered and transcribed in \AMSTeX{} by G. Maltsiniotis. It was edited by M. Künzer, J. Malgoire, and G. Maltsiniotis. The transcription is as faithful as possible to the manuscript. For the few obvious corrections, or the rare comments by the editors, as well as for the original numbering of the manuscript pages, typewriter characters \texttt{[typewriter]} in brackets are used. A question mark in brackets means that one is not sure of the preceding word.

This edition is provisional. Remarks, comments and corrections are welcome. Send a message to:

\textbf{[e-mail address removed]}

\hfill G. Maltsiniotis

\clearpage

% ===================== page 3 =====================
\printedpage{1}
% head: A. Grothendieck : Dérivateurs, Ch. I, Édition des Universités de Montpellier II et Paris VII
% foot: Version 8 juillet 2001

\msp{page 1}

\section*{Derivators}

\subsection*{\textmd{1.}\ Preliminaries on the inverse image functors of presheaves on index categories, and their multiple adjoints.}

$A$ ``category of coefficients'',\\
$I$ ``index category'',\\
$A(I)=\uHom(I^{\circ},A)$ presheaves on $I$ with values in $A$.

If $f : I \to J$, we get:
\[
f^{*} : A(J) \to A(I) \quad ,
\]
and we are interested in the existence of the successive adjoint functors on the left and on the right. Provisional notation:
\[
f^{?}\ f_{!}\ f^{*}\ f_{*}\ f^{!} \quad .
\]

Here are the main facts:

\textit{a}) Subject to the existence of the relevant $\varprojlim$ (resp. $\varinjlim$), $f_{*}$ (resp. $f_{!}$) exists. For example if there exist \texttt{[some]} $\varprojlim$ (resp. $\varinjlim$) that are \textit{finite}, then $f_{*}$ (resp. $f_{!}$) exists if $I$, $J$ are finite categories.

\textit{b}) Even if $A$ is an abelian category, with $\varprojlim$ and $\varinjlim$ (for example $\Ab$), the \textit{biadjoints} do not always exist. For example, in the particular case where $J=e$ is the one-point category, and $I$ the ordered set
\[
\begin{tikzcd}
 & a \\
0 \arrow[ur] \arrow[dr] & \\
 & b
\end{tikzcd}
\]
(in the case of $f^{!}$), or the opposite ordered set (in the case of $f^{?}$). Even if $f$ is an open immersion, $f^{!}$ does not necessarily exist (for example it does not exist if $I$ is as \texttt{[here-]}above, and
\[
J=\begin{tikzcd}
 & a \arrow[dr] & \\
0 \arrow[ur] \arrow[dr] & & c \\
 & b \arrow[ur] &
\end{tikzcd}
\qquad\qquad \text{),}
\]
and if $f$ is a closed immersion, $f^{?}$ does not necessarily exist.

\textit{c}) If $f$ is a \textit{closed immersion}, and $A$ is \textit{pointed} (\textit{i.e.} admits an object both initial and final), then $f^{!}$ exists. If $f$ is an open immersion, $f^{?}$ exists. (Always subject
\printedpage{2}
% head: A. Grothendieck : Dérivateurs, Ch. I, Édition des Universités de Montpellier II et Paris VII
% foot: Version 8 juillet 2001
to the existence in $A$ of the relevant limits). If $f$ is an immersion, $f_{*}$ and $f_{!}$ \texttt{[are]} fully faithful \textit{i.e.} $f^{*}$ \texttt{[is]} a localization functor.

\textit{d}) Let us assume only the existence of the final object $e$ in $A$ (plus if need be certain finite $\varprojlim$). If

\msp{page 2}
\[
i : Y \hookrightarrow S
\]
is a closed immersion, with complement $U$
\[
Y \overset{i}{\hookrightarrow} S \overset{j}{\hookleftarrow} U
\]
so that $j$ is an open immersion, then $A(Y)$ is identified by the functor $i_{*}$ with the full subcategory of $A(S)$ consisting of the $F \in \Ob\ A(S)$ such that $F|U \simeq e_{U}$ (final presheaf). Similarly, if the initial object $\emptyset$ exists in $A$, and the relevant $\varinjlim$, then $A(U)$ is identified by $j_{!}$ with the full subcategory of $A(S)$ consisting of the $F$ such that $F|Y=\emptyset_{Y}$ (initial presheaf, taking the value $\emptyset$ everywhere).

If $A$ is even \textit{pointed} ($e=\emptyset$), then we have, for every $F\in \Ob\ A(S)$, a square both cartesian and cocartesian, functorial in $F$
\[
\begin{tikzcd}
 & F \arrow[dr, "\text{adj}"] & \\
j_{!}j^{*}F \arrow[ur, "\text{adj}"] \arrow[dr] & & i_{*}i^{*}F \\
 & e \arrow[ur] &
\end{tikzcd}
\qquad\qquad \text{.}
\tag*{(1)}
\]
In the case $A$ abelian, this diagram takes the form of an exact sequence
\[
0 \longrightarrow j_{!}j^{*}F \longrightarrow F \longrightarrow i_{*}i^{*}F \longrightarrow 0 \quad .
\]

\textit{e}) Same conditions for $i$ and $j$. We assume that $A$ has the relevant $\varprojlim$, including final object. We then have a cartesian square functorial in $F\in \Ob\ A(S)$ \texttt{[one must assume the category} $A$ \texttt{pointed]}
\[
\begin{tikzcd}
F \arrow[d, "\text{adj}"'] & i_{*}i^{!}F \arrow[l] \arrow[d] \\
j_{*}j^{*}F & e \arrow[l]
\end{tikzcd}
\tag*{(2)}
\]
\textit{i.e.} $i^{!}F$ is identified by the fully faithful functor $i_{*}$ with the \textit{kernel} of $F\to j_{*}j^{*}F$. In the abelian case, this takes the form of an exact sequence
\[
0 \longrightarrow i_{*}i^{!}F \longrightarrow F \longrightarrow j_{*}j^{*}F \quad .
\]

\clearpage

% ===================== page 5 =====================
\printedpage{3}
% head: A. Grothendieck : Dérivateurs, Ch. I, Édition des Universités de Montpellier II et Paris VII
% foot: Version 8 juillet 2001

\msp{page 3}

I dispense with writing down the dual cocartesian square, when $A$ has the relevant $\varinjlim$, including initial object. It seems that it has not been of use in practice. The classically useful adjoint and biadjoint functors are included below (where the ``not useful'' ones are in brackets)
\[
\begin{aligned}
{}[i_{!}]\ i^{*}\ i_{*}\ i^{!} \\
[j^{?}]\ j_{!}\ j^{*}\ j_{*}
\end{aligned}
\]
\textit{i.e.} the functors $i_{!}$ and $j^{?}$ have not been of use classically.

\textit{f}) Case of the categories $\bDelta_{i}$ (standard simplices, $i\in\mathbb{N}$), and of the functors between them (semi-simplicial maps). It would seem that for such an $f$, if $A$ is pointed, there are \textit{bi-adjoints} both on the left and on the right. Moreover, for $\bDelta_{i}$ and $\bDelta_{j}$ fixed when we look at all the $f:\bDelta_{i}\to\bDelta_{j}$ and all the $g:\bDelta_{j}\to\bDelta_{i}$, that the total sequences of multi-adjoint functors for each of the corresponding $f^{*}$, $g^{*}$ are all the same \texttt{[this is not true]}. That looks true at least if $j=i+1$ (and I have checked it for $i=0$, $i=1$), we should find a ``sequence of adjunctions'' of length $2i+7=(i+1)+(i+2)+2+2$. To be precise, the sequence has four ``extreme'' terms (the first two, the last two) and the $(i+1)+(i+2)$ intermediate terms corresponding to the functors of the form $f^{*}$ or $g^{*}$, where $f:\bDelta_{i+1}\to\bDelta_{i}$ is one of the $(i+1)$ degeneracy maps, and $g:\bDelta_{i}\to\bDelta_{i+1}$ one of the $i+2$ face maps \texttt{[this is not quite true; one must restrict oneself to the maps that are injective or surjective]}.